% Discussion
% Interpret results; compare with literature. Keep this in your own voice.

\chapter{Discussion}\label{ch:discussion}

\section{Interpretation of Results}
% TODO: What the latency and engagement results mean.

\section{Comparison with Related Work}
% TODO: How your system compares to prior VR rehab and tracking systems.

\section{Limitations}

Technical and study limitations include the pilot sample size, the scope of outcome measures, and the lack of standardised clinical instruments. Existing work notes a lack of standardised protocols, data formats, and outcome measures in ML-based VR rehabilitation, which complicates direct comparison across systems \cite{wang2025airehab, plavoukou2025sensors}. This pilot used system feel (perceived responsiveness) and engagement as primary outcomes; it did not use clinical patient-reported outcome measures (PROMs) such as KOOS or IKDC. Future work could align with such PROMs and with evidence-based rehabilitation timelines to support clinical translation and validation. Privacy-preserving approaches such as federated learning and edge AI are emerging for personalised rehabilitation and were out of scope for this prototype but are relevant for future at-home deployment \cite{piechowiak2025wearables, chen2025federated}.

\section{Implications for Practice}
% TODO: Implications for clinicians and future systems.
